\chapter{Ringraziamenti}
\section{Buongiorno}
Mi sveglio, é davvero una bella giornata. É domenica, non si va a scuola, posso rimanere tutta la mattina a giocare alla PS2. A pranzo si va dai nonni a mangiare le cotolette e ci saranno tutti. Quando arrivo a casa dei miei nonni mi trovo tutti davanti che mi sorridono. \textbf{Nonno Savio, Nonna Lina, Nonno Filippo, Nonna Gina, Zio Franco, Zia Giovanna, Zia Maria Lucia, Zia Rossana, Zio Aldo, Zia Barbara} stanno li a guardarmi sorpresi, mamma ha appena raccontato loro che ho imparato a contare fino a 100. Non capisco tutto questo stupore, in fondo un bambino di 6 anni bene o male ha giá imparato un po' a leggere, non vedo perché i numeri dovrebbero essere diversi. Tutti mi fanno i complimenti e dicono che devo mettermi a studiare, che ho un futuro davanti, ma a me non piace studiare, molto meglio la PS.
\vspace{5mm}

\textit{Forse in fondo un po' mi piaceva quello che mi dicevano, mi lusingava, ma mi caricava anche di responsaabilitá che forse non avrei voluto. Ciononostante, a saperlo giá 18 anni fa quello che mi sarebbe successo, avrei accettato queste responsabilitá almeno altre 100 volte, tante quanti i numeri che avevo imparato a mettere in ordine. In me hanno visto qualcosa di luminoso, ci hanno creduto che ci fosse, ma in fondo ognuno ha la sua scintilla, che illumina chi ha intorno di un colore o di un altro. Volevo davvero essere perfetto per loro, il piú piccolo della famiglia (almeno al momento), ma per cui tutti avevano le aspettative piú grandi.}
\vspace{5mm}

Torno a casa, stanco morto. Vorrei solo passare il resto della giornata attaccato alla tv, tanto non ho molto altro da fare, ma invece arrivano i miei cugini a casa, \textbf{Stefano, Marco, Andrea, Fabrizio, Alessandra}. Alcuni abitano vicino casa, altri leggermente piú lontano, ma nonostante ció sono venuti tutti per giocare insieme a me con i lego. Furtivi vogliono attingere alla mia preziosa collezione di pezzi particolari, ma non mi faccio abbindolare! Ci sono tutti, vogliono farmi divertire, un po' mi prendono in giro, ma in fondo sono il piú piccolo, un po' me lo merito. Quando saró grande saró sicuro e figo come loro, me lo prometto con tutto me stesso. Hanno le idee cosí chiare sul futuro, e poi sono cosí bravi con i lego, li stimo un sacco!
\vspace{5mm}

\textit{Cosí adulti ai miei occhi, ma senza che lo sapessi cosí insicuri, forse anche loro, consapevoli come adesso anche io un po' lo sono diventato. La stima non ha fatto altro che rafforzarsi, vedendo come quelle insicurezze che ho dovuto passare anche io le hanno affrontate fino a brillare, di quella luce, ognuno a modo suo, cosí realizzati. Siete stati dei modelli, non IL modello, quello spetta a qualcun altro, ma sicuramente quelle persone da cui ho potuto imparare tanti di quegli errori ed evitarli prima che colpissero me stesso. Non avete idea di quanto vi sono grato.}
\vspace{5mm}

A quanto pare hanno regalato un nuovo telescopio ad \textbf{Antonio}, e stasera il cielo é particolarmente pulito, quindi da casa sua possiamo vedere per bene le stelle. Facciamo a turno per sbirciare attraverso quel buco di cielo proiettato sulla lente, le guardo a fondo ma non so esattamente cosa siano, non ho molta idea di quali stelle siano quali, ma in un modo o in un altro mi affascinano. C'é \textbf{Serena} in braccio a zio che ha appena finito di fare la pappa e ora ci fa compagnia in soggiorno. É cosí piccola, ride per qualsiasi faccia buffa le mostri, e un po' fa ridere anche me. É proprio vero che sono contagiosi i sorrisi. Ogni volta che la vedo cosí mi chiedo quando é che riusciró a parlarle in modo sensato, che mi racconterá come si sentiva quando rideva cosí. Alla fine mi sono addormentato sul divano, e senza accorgermene sono riapparso nel letto. Che magia!
\vspace{5mm}

\textit{É sempre stato strano essere da un lato il piú piccolo, e dall'altro il piú grande (anche se solo di 14 giorni). Da questo lato peró mi sentivo al loro pari, che avrei potuto condividere tutta la mia vita con loro e curioso di come i nostri cammini si sarebbero separati, ognuno per per la sua strada. Grazie per essere stati sempre al mio fianco fino all'inizio dell'universitá, e per avermi mostrato com'é la prospettiva quando sei tu il piú grande, quando gli anni in piú ce li hai tu. Anche se ci sentiamo poco, so che ci siete ancora qua vicino a me.}

\separatore

Mi sveglio, é una giornata nuvolosa. Devo andare a scuola, oggi inizia l'ultimo anno di Liceo, devo vedere dove sta la mia nuova classe e ritrovare i miei compagni di classe, anche se non sono loro le persone che sono piú felice di rivedere. Oggi ho dormito nel letto di \textbf{Fil}, é da un po' che ci dormo io invece di lui, ma é anche da un po' che lui non é piú a casa. Mi manca, ma tutti crescono, anche mio fratello. 







\textbf{Mamma} mi sveglia ogni 5 minuti, non ho voglia di alzarmi dal letto, ma sa sempre essere convincente, anche a costo di tirarmi via le coperte. Ha gia preparato tutto il necessario per fare colazione, e la sera prima si é anche messa a cucinare diverse cose, cosí che quando torno potró mangiare un ottimo pranzo. Non ho idea di come faccia ad essere giá sveglia dalle 6:30 ed essere cosí attiva. Anche per lei é il primo giorno di scuola, peró sembra molto piú contenta di me di andarci, e ci credo, lei ci insegna, mica deve studiare, quello é pesante (mai cosa piú falsa fu pensata)
\vspace{5mm}

\textit{Mamma, non so come fai, hai dovuto portare avanti uno dei lavori piú importanti ma meno soddisfacenti del mondo, e nel frattempo sei stata anche madre mia e di mio fratello. Instancabile, Incrollabile, Indistruttibile, sempre li quando avevo bisogno con tutta la discretezza, a volte senza che neanche me ne accorgessi. Sempre davanti a sforzi enormi, affrontandoli come se fossero un pacco di noccioline da sgusciare, come quelle che ti piacciono tanto. Vorrei dire che ti sono profondamente grato e che ti stimo piú di chiunque altro, ma queste parole non varrebbero niente comparate a all'indescrivibile rilevanza che hai avuto nella mia crescita.}

\section{Buonasera}

\section{Buonanotte}